\documentclass{article}
\usepackage[utf8]{inputenc}

\usepackage[a4paper, hmargin=3cm, top=3cm, bottom=3cm]{geometry}

%\usepackage{euler} %usa la fuente "euler" para las ecuaciones
\usepackage{amsfonts} %para Z estilizada
\usepackage{pifont} %para más estilos de viñetas
\usepackage{amsmath} %necesario para equation*
\usepackage{amsthm} %necesario para proof enviroment
\usepackage{dsfont} %alguna letra caligráfica
\usepackage{tipa} %Epsilon bonito
\usepackage{graphicx} %para manejar imágenes
\usepackage{wrapfig} %para imágenes en modo "wrap"
\usepackage{lipsum} %para usar el texto de relleno
\usepackage{tikz-cd} %para los diagramas conmutativos
\usepackage{amssymb}
\usepackage{tabularx}
\graphicspath{ {./Images/} } %la dirección de la cual se sacan las imágenes
\usepackage{euler}
\usepackage[normalem]{ulem} % Subrayado
\usepackage{setspace} %Para interlineado

\setlength{\parindent}{0pt}

\renewcommand{\figurename}{Figura}

\title{
  \vspace{-1.5cm}
  \textsc{
    \LARGE{Lógica Matemática\\ Primer parcial}\\ \vspace{1cm}
    \large{Abril 21 de 2026\\ \vspace{1cm} {Presentado por:} \\ \vspace{0.3cm} {Juan Camilo Lozano Suárez}}\\
  }
}
\date{}

\begin{document}
    
    \maketitle
    \subsection*{\uline{Ejercicio1}}
Sea $\mathcal{E} = \{ \forall x \exists y P(x,y), \forall x \forall y (P(x,y) \rightarrow Q(y)), \forall x (Q(x) \rightarrow R(x)) \}$. Pruebe que $\mathcal{E} \vdash \forall x \exists y R(y)$.\\

\textbf{Solución:} Asumimos que $\mathcal{E}$ contiene los axiomas del lenguaje. Llamamos $E_1: \forall x \exists y P(x,y)$, $E_2: \forall x \forall y (P(x,y)\rightarrow Q(y))$ y $E_3: \forall x (Q(x) \rightarrow R(x))$. La siguiente es una deducción de $\forall x \exists y R(y)$ a partir de $\mathcal{E}$:

\begin{enumerate}
    \item $\forall x(Q(x) \rightarrow R(x))$ \hfill(Premisa $E_3$)
    \item $\forall x(Q(x) \rightarrow R(x)) \rightarrow (Q(y) \rightarrow R(y))$ \hfill(Axioma de especialización, $t = y$)
    \item $Q(y) \rightarrow R(y)$ \hfill(MP 1, 2)
    \item $\forall x \forall y(P(x, y) \rightarrow Q(y))$ \hfill(Premisa $E_2$)
    \item $\forall x \forall y(P(x, y) \rightarrow Q(y)) \rightarrow \forall y(P(x, y) \rightarrow Q(y))$ \hfill(Axioma de especialización, $t = x$)
    \item $\forall y(P(x, y) \rightarrow Q(y))$ \hfill(MP 4, 5)
    \item $\forall y(P(x, y) \rightarrow Q(y)) \rightarrow (P(x, y) \rightarrow Q(y))$ \hfill(Axioma de especialización, $t = y$)
    \item $P(x, y) \rightarrow Q(y)$ \hfill(MP 6, 7)
    \item $(P(x, y) \rightarrow Q(y)) \rightarrow ((Q(y) \rightarrow R(y)) \rightarrow (P(x, y) \rightarrow R(y)))$ \hfill(Tautología: Transitividad)
    \item $(Q(y) \rightarrow R(y)) \rightarrow (P(x, y) \rightarrow R(y))$ \hfill(MP 8, 9)
    \item $P(x, y) \rightarrow R(y)$ \hfill(MP 3, 10)
    \item $(P(x, y) \rightarrow R(y)) \rightarrow (\neg R(y) \rightarrow \neg P(x, y))$ \hfill(Tautología: Contraposición)
    \item $\neg R(y) \rightarrow \neg P(x, y)$ \hfill(MP 11, 12)
    \item $\forall y (\neg R(y) \rightarrow \neg P(x, y))$ \hfill(Gen. sobre $y$ en 13)
    
    % --- DESGLOSE DEL PASO DE DISTRIBUCIÓN ---
    \item $\forall y \neg R(y)$ \hfill(Suposición provisional para el Lema de la Deducción)
    \item $\forall y (\neg R(y) \rightarrow \neg P(x, y)) \rightarrow (\neg R(y) \rightarrow \neg P(x, y))$ \hfill(Axioma de especialización, $t=y$)
    \item $\neg R(y) \rightarrow \neg P(x, y)$ \hfill(MP 14, 16)
    \item $\forall y \neg R(y) \rightarrow \neg R(y)$ \hfill(Axioma de especialización, $t=y$)
    \item $\neg R(y)$ \hfill(MP 15, 18)
    \item $\neg P(x, y)$ \hfill(MP 17, 19)
    \item $\forall y \neg P(x, y)$ \hfill(Gen. sobre $y$ en 20)
    \item $\forall y \neg R(y) \rightarrow \forall y \neg P(x, y)$ \hfill(Lema de la deducción aplicado a 15-21)
    % -----------------------------------------

    \item $(\forall y \neg R(y) \rightarrow \forall y \neg P(x, y)) \rightarrow (\neg \forall y \neg P(x, y) \rightarrow \neg \forall y \neg R(y))$ \hfill(Tautología: Contraposición)
    \item $\neg \forall y \neg P(x, y) \rightarrow \neg \forall y \neg R(y)$ \hfill(MP 22, 23)
    \item $\exists y P(x, y) \rightarrow \exists y R(y)$ \hfill(Definición de $\exists$ y eliminación de doble negación)
    \item $\forall x \exists y P(x, y)$ \hfill(Premisa $E_1$)
    \item $\forall x \exists y P(x, y) \rightarrow \exists y P(x, y)$ \hfill(Axioma de especialización, $t = x$)
    \item $\exists y P(x, y)$ \hfill(MP 26, 27)
    \item $\exists y R(y)$ \hfill(MP 25, 28)
    \item $\forall x \exists y R(y)$ \hfill(Gen. sobre $x$ en 29)
\end{enumerate}

\hfill \qed

    \subsection*{\uline{Ejercicio2}}
Sea $\mathcal{L} = \{ <, 0 \}$ un lenguaje con una relación binaria $<$ y una constante $0$. Considere la teoría $T$ con los axiomas:
    \begin{itemize}
        \item $\forall x \neg(x < 0)$
        \item $\forall x \forall y \forall z ((x < y \wedge y < z) \rightarrow x < z)$
        \item $\forall x \exists y (x < y)$
    \end{itemize}
    Sea $c$ un nuevo símbolo de constante. Muestre que $T_c = T \cup \{ \neg(c < 0) \}$ es consistente. \\
    Muestre que $T$ tiene un modelo contable y construya explícitamente un modelo $\mathcal{M} = (M, <, 0)$ que satisfaga los axiomas de $T$ y contenga un elemento $c$ tal que la interpretación de $c$ no es igual a $0$.\\

\textbf{Solución:}


    \subsection*{\uline{Ejercicio3}}
Sea $\mathcal{E} = \{ \forall x (\forall y (P(y) \rightarrow Q(x)) \rightarrow R(x)), \forall x (\exists y P(y) \rightarrow S(x)), \forall x (S(x) \rightarrow \exists y T(y)), \exists x \neg R(x), \forall y \neg T(y) \}$. \\
    Pruebe que $\mathcal{E} \vdash \exists x \neg Q(x)$.\\

\textbf{Solución:} Asumimos que $\mathcal{E}$ contiene los axiomas del lenguaje. Llamamos $E_1: \forall x ( \forall y ( P(y) \rightarrow Q(x) ) \rightarrow R(x) )$, $E_2: \forall x ( \exists y P(y) \rightarrow S(x) )$, $E_3: \forall x ( S(x) \rightarrow \exists y T(y) )$, $E_4: \exists x \neg R(x)$ y $E_5: \forall y \neg T(y)$. Para probar que $\mathcal{E} \vdash \exists x \neg Q(x)$ utilizaremos reducción al absurdo

\begin{enumerate}
   \item $\neg \exists x \neg Q(x)$ \hfill (Nueva premisa)
   \item $\forall x Q(x)$ \quad \hfill (Axioma lógico para cuantificadores y eliminación de doble negación vía tautología)
   \item $\forall y \neg T(y)$ \hfill (Premisa $E_5$)
   \item $\neg \exists y T(y)$ \hfill (Axioma lógico para cuantificadores)
   \item $\forall x ( S(x) \rightarrow \exists y T(y) )$ \hfill (Premisa $E_3$)
   \item $S(x) \rightarrow \exists y T(y)$ \hfill (Axioma de especialización, $t=x$)
   \item $\neg \exists y T(y) \rightarrow \neg S(x)$ \hfill (Tautología: Contraposición)
   \item $\neg S(x)$ \hfill (MP 4, 7)
   \item $\forall x ( \exists y P(y) \rightarrow S(x) )$ \hfill (Premisa $E_2$)
   \item $\exists y P(y) \rightarrow S(x)$ \hfill (Axioma de especialización, $t=x$)
   \item $\neg S(x) \rightarrow \neg \exists y P(y)$ \hfill (Tautología: Contraposición)
   \item $\neg \exists y P(y)$ \hfill (MP 8, 10)
   \item $\forall y \neg P(y)$ \hfill (Axioma lógico para cuantificadores)
   \item $\neg P(y)$ \hfill (Axioma de especialización, $t=y$)
   \item $\neg P(y) \rightarrow ( P(y) \rightarrow Q(x) )$ \hfill (Tautología)
   \item $P(y) \rightarrow Q(x)$ \hfill (MP 14, 15)
   \item $\forall y (P(y) \rightarrow Q(x))$ \hfill (Generalización sobre $y$)
   \item $\forall x ( \forall y ( P(y) \rightarrow Q(x) ) \rightarrow R(x) )$ \hfill (Premisa $E_1$)
   \item $\forall y ( P(y) \rightarrow Q(x) ) \rightarrow R(x)$ \hfill (Axioma de especialización, $t=x$)
   \item $R(x)$ \hfill (MP 17, 19)
   \item $\forall x R(x)$ \hfill (Generalización sobre $x$)
   \item $\neg \exists x \neg R(x)$ \hfill (Axioma lógico para cuantificadores y uso de la tautología de doble negación)
   \item $\exists x \neg R(x)$ \hfill (Premisa $E_4$)
\end{enumerate}
Los pasos 22 y 23 nos dan la contradicción deseada. Como la suposición inicial $\neg \exists x \neg Q(x)$ nos lleva a una contradicción con la premisa $E_4$, concluimos que la suposición es falsa, y por lo tanto necesariamente se tiene $\exists x \neg Q(x)$. Esto prueba $\mathcal{E}\vdash \exists x \neg Q(x)$.\\
\hfill \qed

    \subsection*{\uline{Ejercicio4}}
Sea $\mathbb{Z}_2$ el cuerpo de los enteros módulo 2. Exprese los 5 conectivos básicos mediante las operaciones de $\mathbb{Z}_2$.\\

\textbf{Solución:}

    \subsection*{\uline{Ejercicio5}}
Sea $K$ un cuerpo y $p(X_1, \dots, X_n) \in K[X_1, \dots, X_n]$. Sea $\overline{p}$ la función polinomial asociada a $p$, definida por:
    \begin{equation*}
        \overline{p} : K^n \rightarrow K, \quad (a_1, \dots, a_n) \mapsto p(a_1, \dots, a_n)
    \end{equation*}

    Muestre que para toda fórmula proposicional $\alpha(p_1, \dots, p_n)$, existe un polinomio $p_\alpha \in \mathbb{Z}_2[X_1, \dots, X_n]$ tal que para toda valuación $v$ se tenga $\overline{v}(\alpha) = \overline{p}_\alpha(v(p_1), \dots, v(p_n))$. \\

    ¿Vale la unicidad de $p_\alpha$ para una fórmula $\alpha$ dada?\\

\textbf{Solución:}


\end{document}
